\chapter{专用平台开发}

\section{领域导言与演进背景}

\begin{preamblebox}[知识领域概述与变革]
专用平台开发 (SPD)知识领域旨在培养学生针对特定行业场景、特定硬件体系与严苛业务需求，构建专用软件平台与行业解决方案的核心能力。在现代数字化社会中，专用平台已全面渗透至医疗健康、金融科技、交通运输、智能制造、具身智能及元宇宙交互等诸多关键垂直领域。专用平台的工程构建涵盖了从系统需求分析、体系架构设计、异构软硬件协同开发，到全生命周期自动化部署与高可用维护的完整链路。

\textbf{自 CS2013 以来的核心变革：}
\begin{enumerate}
    \item \textbf{新兴计算前沿与垂直领域的爆发式增长}：数据科学与大数据分析依赖多端跨平台感知数据采集；网络空间安全要求在软硬件底层强化数据隔离与信道防御；人工智能与机器学习催生了具身机器人与无人机的自主协同控制；交互式计算平台迎来了大语言模型（LLMs，如 OpenAI ChatGPT、Codex 与 GitHub Copilot）的革命性融合，极大地拓展了代码生成、交互式探索与创意计算的技术边界。
    \item \textbf{产业界工程师核心能力诉求剧增}：行业对精通移动端、现代 Web、机器人、嵌入式硬件及游戏引擎的平台级全栈工程师呈现迫切需求。包括：在碎片化设备与操作系统间打造体验高度一致的移动应用、面向头戴式显示器（HMD）与混合现实（MR）评估无障碍交互体验、基于云原生与现代全栈语言构建高安全 Web 业务底盘，以及利用微控制器（Raspberry Pi、Arduino）和开源机器人操作系统（ROS）实现软硬件互联控制。
    \item \textbf{知识领域全面升级与重命名}：由 CS2013 中的“基于平台的开发 (PBD)”正式重命名为“专用平台开发 (SPD)”，精准反映了平台开发从通用中间件走向垂直专用软硬件生态的时代转型。
    \item \textbf{计算机科学核心学时（计算机核心）实现历史性突破}：在 CS2013 中本领域核心学时为 0；而在 CS2023 中，计算机科学核心学时大幅提升至 \textbf{8$\sim$9 学时}。大纲将 Web 与移动端划分为通识核心基础与专用平台进阶模块，为全球各高校依据自身办学定位灵活配置特色专业方向（如游戏开发、智能硬件、云计算全栈）提供了强大的弹性支撑。
    \item \textbf{吸纳前沿知识单元并强化社会伦理责任}：全面设立“机器人平台”、“嵌入式平台”、“游戏平台”、“交互式计算与大模型平台”，并将“社会、伦理与专业素养（SEP）”深度融入硬件适老化、无障碍设计及端侧数据隐私合规的全流程中。
\end{enumerate}
\end{preamblebox}

\section{核心学时分配建议}

本领域在课程体系中采用“核心基石 + 弹性专精（领域核心/选修）”的双轨学时机制。Web 与移动开发的核心基础构成了全体计算机专业毕业生的通用必修基底（计算机核心 累计 8$\sim$9 学时），而机器人、嵌入式、游戏引擎与交互计算平台则作为各高校计算机专业深化方向的核心选修（领域核心 累计 15 学时）：

\begin{table}[H]
\centering
\setlength{\tabcolsep}{3.5pt}\fontsize{8pt}{9.6pt}\selectfont\begin{tabularx}{\columnwidth}{Y c c}
\toprule
\textbf{知识单元} & \textbf{计算机核心} & \textbf{领域核心} \\
\midrule
跨平台共性技术与关注点 & 2 & -- \\
Web 平台全栈开发与云原生架构 & 3 & -- \\
移动计算平台与端侧智能应用 & 3 & -- \\
机器人系统与具身智能开发平台 & -- & 3 \\
嵌入式系统与低功耗物联网平台 & -- & 3 \\
游戏引擎架构与实时交互互动平台 & -- & 4 \\
交互式计算平台与大模型应用 & -- & 2 \\
平台社会责任、伦理合规与专业素养 & -- & 3 \\
\midrule
\textbf{合计} & \textbf{8 学时} & \textbf{15 学时} \\
\bottomrule
\end{tabularx}
\end{table}

\noindent\textbf{【学时弹性说明】}：部分办学方案中可将 Web 基础（2 学时）与移动端基础（1 学时）抽离为先导基础（计算机核心 共 3 学时），再结合高校特色在各专用平台（*）中选配 6 学时深度实践，总计覆盖 9 门类核心计算机学时。

\section{知识单元与详细大纲}

\subsection{跨平台共性技术与架构关注点}
\textbf{【学时建议】}：2 计算机核心学时（计算机核心）。

本单元旨在系统梳理面向多样化专用计算平台进行软件构建时的通用方法论与核心架构约束。引导学生深入理解不同平台（移动端、Web、游戏主机、工业自动化、嵌入式设备）的软硬件异构特性，掌握专用应用编程接口（API）与跨平台编程语言的高效应用方法。

\textbf{【核心知识点与专题大纲】}
\begin{itemize}
    \item \textbf{现代专用平台生态全景概览}：Web 浏览器平台、移动终端平台、主机与掌机游戏平台、工业自动化与机器人系统平台。
    \begin{itemize}
        \item \textbf{输入与感知控制机制}：多点触控、空间惯性传感器（IMU）、计算机视觉跟踪（Inside-out/Outside-in）、触觉力反馈。
        \item \textbf{严苛平台资源约束分析}：有限算力与热耗限制；受限板载内存与持久化存储容量；非对称与弱网无线通信带宽；系统高可用性、容错韧性与工业合规安全性。
        \item \textbf{输出与受控执行机构}：高刷新率与动态色域显示引擎、物理执行机构、高精度力矩伺服电机与震动反馈装置。
    \end{itemize}
    \item \textbf{专用平台 API 驱动式编程 vs. 传统通用软件构造}：硬件抽象层（HAL）与操作系统内核沙箱的交互契约；平台级生命周期事件驱动模型 vs. 传统主控制台执行流。
    \item \textbf{跨平台现代开发语言图谱与选型}：Kotlin、Swift、C\#、C++、Java、JavaScript、TypeScript 以及 HTML5 Web 标准。
    \item \textbf{平台约束环境下的工程构建}：受限交叉编译工具链、真机远程联合调试、内存基线分析与安装包体积削减策略。
    \item \textbf{专用领域语言与新兴平台的快速掌握范式}：语言运行时特性的迁移学习、平台官方架构指南研读法、SDK 核心范式的逆向抽象。
\end{itemize}

\textbf{【示范性学习成果与能力指标】}
\begin{itemize}
    \item 列举并深度剖析移动端开发在内存占用、并发能耗与交互时延方面的核心约束。
    \item 阐明 Web 开发经典三层架构（表示层、业务逻辑层、持久化数据层）在现代富客户端形态下的演化机制。
    \item 解释在 Web 无状态通信协议（HTTP）下利用 Cookie、Session、Token（JWT）及分布式缓存维护用户会话状态的底层原理。
    \item 辨析动态脚本语言与强类型编译型语言（C++/Rust）在运行效率、内存安全及跨平台部署维度的核心权衡。
\end{itemize}

\subsection{平台全栈开发与云原生架构}
\textbf{【学时建议】}：3 计算机核心学时（计算机核心）。

本单元聚焦于现代 Web 生态的体系结构、全栈编程技术、安全防御体系及云原生服务治理。涵盖富客户端现代框架、SaaS 架构、无障碍标准、端到端数据管理以及弹性伸缩机制。

\textbf{【核心知识点与专题大纲】}
\begin{itemize}
    \item \textbf{现代 Web 全栈开发语言体系}：HTML5 语义化标准、CSS3 弹性布局与网格排版系统、现代 ECMAScript（JavaScript）、TypeScript 静态类型系统及后端开发语言（Node.js、PHP、Go、Python）。
    \item \textbf{Web 应用平台与现代前端元框架}：组件化虚拟 DOM 体系、服务端渲染（SSR）、静态站点生成（SSG）、增量静态再生（ISR）以及全栈集成框架。
    \item \textbf{软件即服务设计范式}：多租户架构隔离（数据隔离与模式隔离）、订阅式计量与权限网关。
    \item \textbf{Web 国际标准与无障碍体验}：W3C 文档对象模型（DOM）规范、无障碍网页内容指南（WCAG 2.1）、语义化 ARIA 标签规范。
    \item \textbf{Web 平台全栈安全与隐私纵深防御}：跨站脚本攻击（XSS）、跨站请求伪造（CSRF）、SQL 注入防御；内容安全策略（CSP）、HTTP 严格传输安全（HSTS）、CORS 跨域资源共享策略。
    \item \textbf{Web 业务系统需求分析与架构权衡}：用户高并发峰值预测、首屏加载时间（FPT）、累计布局偏移（CLS）与全链路核心体验指标。
    \item \textbf{现代云端算力基础设施与服务治理}：主流公有云（AWS、Azure、阿里云等）计算生态。
    \begin{itemize}
        \item \textbf{云端托管与容器化}：虚拟机、Docker 容器化与 Kubernetes 编排集群。
        \item \textbf{系统弹性与伸缩}：自动伸缩组、多可用区负载均衡（ALB/SLB）、CDN 边缘智能缓存。
        \item \textbf{云服务容量规划与成本工程}：按需计费、预留实例分析及 FinOps 成本优化模型。
    \end{itemize}
    \item \textbf{全栈数据生命周期治理}：
    \begin{itemize}
        \item \textbf{数据驻留与属地合规}：跨国数据合规路径、跨境数据流动加密与主权审计。
        \item \textbf{数据完整性与合规销毁}：数据多副本强一致性容灾机制、用户注销后的彻底可信擦除。
    \end{itemize}
    \item \textbf{现代软件架构范式演进}：
    \begin{itemize}
        \item \textbf{单体架构 vs. 微服务架构}：服务解耦、领域驱动设计（DDD）、分布式事务（Saga 模式）。
        \item \textbf{微前端架构}：独立构建与部署、模块联邦、运行时沙箱隔离。
        \item \textbf{事件驱动架构（EDA） vs. RESTful API 架构}：异步消息队列高吞吐解耦 vs. 同步轻量契约通信的优劣权衡。
        \item \textbf{无服务器计算}：函数即服务、冷启动优化、事件源触发与按需弹性计费。
    \end{itemize}
    \item \textbf{持久化存储方案综合选型}：
    \begin{itemize}
        \item \textbf{关系型数据库（RDBMS）}：ACID 事务、索引 B+ 树优化、读写分离与主从复制。
        \item \textbf{NoSQL 与 NewSQL 体系}：文档型数据库、键值数据库（Redis 缓存防穿透/击穿/雪崩）、分布式列存数据库。
    \end{itemize}
\end{itemize}

\textbf{【示范性学习成果与能力指标】}
\begin{itemize}
    \item 遵循微服务或无服务器架构规范，独立设计并实现一套具备弹性伸缩能力的 Web 业务原型系统。
    \item 全面剖析 Web 运行平台带来的机遇与瓶颈，重点阐述容器托管、动态扩缩容与跨域跨端安全通信的实现路径。
    \item 对比分析 Web 敏捷开发范式与传统单机通用软件开发的异同点（在发布频率、状态持久化与客户端异构性层面）。
    \item 阐明 SaaS 架构在持续交付、多租户计费及弹性容量治理上相较于传统买断式桌面软件的核心优势。
    \item 针对现有线上大型 Web 系统执行 WCAG 2.1 无障碍合规性与 W3C 标准审计，提出优化整改技术方案。
\end{itemize}

\subsection{移动计算平台与端侧智能应用}
\textbf{【学时建议】}：3 计算机核心学时（计算机核心）。

本单元专注于移动计算终端（iOS/Android/鸿蒙）的开发技术、底层运行约束及新型交互能力。深入探讨移动端本地端侧计算、专用 NPU 神经网络加速、硬件沙箱安全隔离以及高并发异步流式响应。

\textbf{【核心知识点与专题大纲】}
\begin{itemize}
    \item \textbf{现代移动工程研发技术栈}：
    \begin{itemize}
        \item \textbf{主流移动开发语言}：Swift（现代函数式与强类型特性）、Kotlin（协程高并发与空安全）、Dart。
        \item \textbf{官方集成开发环境}：Xcode、Android Studio 及其真机性能调优工具链。
        \item \textbf{原生开发 vs. 跨平台技术框架}：Native 深度控制力 vs. 跨平台框架（Flutter 统一渲染管道、React Native 桥接架构）的开发效率与运行开销权衡。
        \item \textbf{移动端主流架构模式}：MVC、MVP、MVVM（数据绑定）与现代单向数据流架构（MVI）。
    \end{itemize}
    \item \textbf{移动平台硬性软硬件约束与优化}：
    \begin{itemize}
        \item \textbf{用户界面与触控交互设计}：响应式屏幕适配、沉浸式手势导航、深色模式与动态字体无障碍支持。
        \item \textbf{移动 App 与 Web 页面用户体验差异}：离线可用性、平滑转场动画（60/120 FPS 无卡顿刷新）与微交互设计。
        \item \textbf{端侧系统安全防护}：应用进程沙箱隔离、运行时动态权限最小化申请、安全飞地与 Keyring/Keystore 加密密钥存储。
        \item \textbf{功耗与计算性能权衡}：电量消耗模型、后台唤醒频率对齐、动态降频应对发热降频。
    \end{itemize}
    \item \textbf{移动数据互联与轻量接口设计}：
    \begin{itemize}
        \item \textbf{移动专属 API 端点优化}：针对弱网（高时延、高丢包）的批量请求合并、增量数据同步、BFF（Backend For Frontend）架构。
        \item \textbf{本地缓存与离线数据同步}：本地轻量 SQLite/Room/CoreData 存储与断网重连自动补偿机制。
    \end{itemize}
    \item \textbf{移动计算特有使能技术}：
    \begin{itemize}
        \item \textbf{空间位置感知计算}：高精度 GNSS、基站三角定位、Wi-Fi 室内指纹与地理围栏低功耗监测。
        \item \textbf{多模态传感器融合计算}：三轴陀螺仪、加速度计、气压计、健康体征传感器（光学心率/血氧）数据采集与异常过滤。
        \item \textbf{实时通讯与系统级通道}：移动电话通信、VoIP 低时延流媒体、系统级推送机制。
        \item \textbf{移动端增强现实（AR）}：基于 ARKit/ARCore 的平面检测、六自由度（6DoF）空间定位与虚拟物体虚实遮挡融合。
    \end{itemize}
    \item \textbf{工程规范、机型适配与自动化测试}：
    \begin{itemize}
        \item 碎片化屏幕与多版本 OS 兼容性云测、自动化 UI 测试、端侧灰度发布与崩溃埋点上报。
    \end{itemize}
    \item \textbf{移动端异步计算与响应式编程}：
    \begin{itemize}
        \item 异步非阻塞模型与同步阻塞模型的底层区别；主线程（UI 线程）避免掉帧（ANR 规避）。
        \item 现代异步控制范式：回调机制、Promise、async/await 以及结构化并发（Kotlin 协程 Coroutines / Swift Concurrency Tasks）。
        \item 异步错误捕获传播链与网络超时异常降级策略。
        \item 异步并发代码测试难题：竞态条件排查、时间依赖测试与内存悬空引用检测。
    \end{itemize}
\end{itemize}

\textbf{【示范性学习成果与能力指标】}
\begin{itemize}
    \item 实现一款基于位置感知（LBS）的移动应用，能够利用后台地理围栏与云端数据 API 实现动态信息实时推送。
    \item 构建一款传感器驱动的高性能移动应用，高效采集陀螺仪与体征数据并以低功耗模式同步至云端服务器。
    \item 结合原生通信能力与 WebRTC/WebSocket 协议，编写一款支持即时文本与低时延音视频通话的移动应用。
    \item 辨析原生开发与基于 Flutter/React Native 的跨平台技术方案，在首屏冷启动速度、动画渲染性能及团队维护成本上的核心利弊。
\end{itemize}

\subsection{机器人系统与具身智能开发平台}
\textbf{【学时建议】}：3 领域核心学时（领域核心）。

本单元探讨针对现有实体与仿真机器人平台进行系统级软件构建与具身智能部署的技术。内容涵盖机器人软硬件架构、专用开发环境与库（如机器人操作系统 ROS/ROS2）、空间运动学与动力学、自主导航与路径规划，以及物理世界人机协作安全规范。

\textbf{【核心知识点与专题大纲】}
\begin{itemize}
    \item \textbf{现代机器人平台与设备类型}：轮式移动机器人（AGV/AMR）、工业多轴机械臂、足式四足/双足人形机器人、无人飞行器（UAV）。
    \item \textbf{感知、嵌入式板载计算与受控执行器}：激光雷达、RGB-D 深度相机、超声波测距传感器；无刷伺服电机、高精度减速机（谐波/RV）与末端电磁抓夹。
    \item \textbf{机器人专用开发语言、框架与生态库}：C++ / Python 系统级接口、URDF/Xacro 机器人物理模型描述文件、MoveIt 运动规划框架、OpenCV 视觉与 PCL 点云处理库。
    \item \textbf{机器人平台物理约束与工程考量}：高精度实时时序约束、非完整约束、动力矩饱和限制与电池续航极限。
    \item \textbf{虚实融合交互与仿真引擎互联}：基于 Gazebo、Webots 或 NVIDIA Isaac Sim 建立高保真物理仿真模型并实现零成本迁移。
    \item \textbf{机器人学核心基础算法}：
    \begin{itemize}
        \item \textbf{软件体系架构}：机器人操作系统 ROS/ROS2 核心通信机制（节点 Node、话题 Topic 异步发布/订阅、服务 Service 同步请求、动作 Action 目标打断反馈机制）。
        \item \textbf{正向运动学}：空间位姿描述、齐次变换矩阵、DH 参数法。
        \item \textbf{逆运动学}：几何解析法与数值迭代法（雅可比矩阵求逆与阻尼最小二乘法）。
        \item \textbf{刚体动力学}：拉格朗日力学与牛顿-欧拉迭代动力学方程、转动惯量张量计算。
        \item \textbf{自主导航与路径规划}：激光/视觉 SLAM 同步定位与建图、全局规划（$A^*$ 算法、Dijkstra、RRT*）与局部避障控制（TEB/DWA 算法）。
        \item \textbf{末端操控与抓取规划}：空间轨迹插补（样条曲线）、末端六自由度抓取姿态估计与力控自适应抓取。
        \item \textbf{人机协作物理安全规范}：急停联锁保护机制、碰撞力矩检测顺应控制与 ISO 10218/TS 15066 安全标准。
    \end{itemize}
\end{itemize}

\textbf{【示范性学习成果与能力指标】}
\begin{itemize}
    \item 在指定机器人软硬件平台（如 ROS 连接的实体移动底盘或高保真物理仿真器）上独立设计并实现一套自主导航任务流。
    \item 基于微控制器（如 Arduino/树莓派）组装一套双轮差速寻迹机器人，编写底层嵌入式控制代码使其自主穿越未知迷宫。
    \item 结合计算图架构，深入阐释 ROS 发布/订阅机制与通用多线程操作系统进程间通信（IPC）在时间确定性与资源开销上的异同。
    \item 针对给定的工业物流搬运任务，使用 ROS2 架构设计一套完整的软硬件协同系统，明确定义各功能节点及其 Topic 与 Action 接口规范。
\end{itemize}

\subsection{嵌入式系统与低功耗物联网平台}
\textbf{【学时建议】}：3 领域核心学时（领域核心）。

本单元专注于资源严重受限的嵌入式微处理器、片上系统以及广泛泛在计算的物联网边缘节点系统开发。涵盖硬实时确定性、外设总线驱动、超低功耗休眠管理与软硬件协同测试验证。（并发、死锁与底层调度等通用议题可参照 PDC 与 OS 领域）。

\textbf{【核心知识点与专题大纲】}
\begin{itemize}
    \item \textbf{嵌入式系统的本质特征与分类}：
    \begin{itemize}
        \item 硬实时、软实时与非实时系统的时延确定性界定。
        \item 严苛资源受限环境：有限 RAM 堆栈空间、无虚拟内存 MMU 支持、严格硬性截止时间。
    \end{itemize}
    \item \textbf{系统功能安全分析与可靠性设计}：故障安全与故障容错架构；危险与可操作性分析（HAZOP）、故障树分析（FTA）及安全编码规范（如 MISRA C/C++）。
    \item \textbf{物理感知与执行外设硬件接口}：通用输入输出（GPIO）、模拟-数字转换（ADC/DAC）、脉宽调制（PWM）；常见工业总线协议（UART、SPI、I2C、CAN 总线及 Modbus）。
    \item \textbf{嵌入式核心开发范式}：裸机前后台轮询架构、实时操作系统内核（FreeRTOS、RT-Thread、Zephyr RTOS 任务调度）；交叉编译体系与链接器脚本内存分区映射。
    \item \textbf{实时资源调度与并发管理}：速率单调调度（RMS）、最早截止时间优先（EDF）；优先级反转风险与优先级天花板/优先级继承协议（PIP/PCP）；中断服务例程（ISR）短小快设计原则与延迟处理机制。
    \item \textbf{嵌入式系统验证与底层调试技术}：硬件在环仿真（HIL）、静态代码分析、JTAG/SWD 硬件在线断点调试、逻辑分析仪与示波器总线时序抓取。
    \item \textbf{低功耗物联网系统工程设计}：微安级深度休眠唤醒模式、看门狗定时器自动故障自愈、边缘轻量智能计算（TinyML）推理。
\end{itemize}

\textbf{【示范性学习成果与能力指标】}
\begin{itemize}
    \item 针对特定硬件微控制器平台（如 STM32、ESP32），独立设计并开发一套具备传感器采样与低功耗休眠唤醒功能的物联网节点应用。
    \item 深入对比嵌入式实时操作系统（RTOS）与通用分时操作系统（如 Linux）在任务上下文切换、中断响应延迟及内存管理策略上的根本差异。
    \item 编写底层外设驱动代码，通过 SPI/I2C 总线实现对惯性测量单元或温湿度传感器的数据采集，并驱动 OLED 显示。
    \item 熟练运用硬件调试器与逻辑分析仪排查嵌入式系统通信总线偶发数据丢失、时钟毛刺或栈溢出导致的 HardFault 死锁故障。
    \item 结合算力、能耗与物料成本（BOM），评估微控制器（MCU）、数字信号处理器（DSP）与现场可编程门阵列（FPGA）在特定嵌入式应用中的选型合理性。
\end{itemize}

\subsection{游戏引擎架构与实时交互互动平台}
\textbf{【学时建议】}：4 领域核心学时（领域核心）。

本单元专注于在严苛硬件算力约束下构建高性能实时交互图形渲染与互动游戏系统的核心软件工程原理。全面覆盖现代 GPU 硬件渲染管线、物理仿真、着色器编程、主流游戏引擎架构、跨平台移植与人机工效无障碍体验。

\textbf{【核心知识点与专题大纲】}
\begin{itemize}
    \item \textbf{游戏硬件与软件计算平台的演进谱系}：
    \begin{itemize}
        \item \textbf{平台演进历程}：从早期的 Brown Box、家用主机时代走向现代元宇宙、云游戏与数字孪生；CPU 与 GPU 协同异构计算的飞速跃迁；平台融合与移动化趋势。
        \item \textbf{典型游戏平台形态}：高性能个人电脑（PC）、家用主机（PlayStation、Xbox、Switch）、便携掌机、街机、智能电视大屏、移动智能手机、平板电脑、VR/AR 空间计算一体机、沉浸式互动环幕（CAVE 洞穴系统）、Web 浏览器、云端流式渲染平台。
        \item \textbf{不同游戏平台的底层硬件特性与严苛约束}：内存带宽与拓扑结构（DDR vs. GDDR 高带宽高延迟统一内存架构）；GPU/CPU 时钟频率与核心集群调度；专用协处理器（物理加速卡 PPU、向量协处理器 SIMD）；总线带宽（PCIe 通道数）；网络吞吐与往返时延（RTT）；最大函数调用栈深度限制；发热温控降频机制与整机功耗预算；字节序（大端/小端 Endianness）差异。
        \item \textbf{典型交互传感器、控制器与触觉执行机构}：键鼠、游戏手柄摇杆与线性扳机、空间体感可穿戴设备、触控交互表面；计算机视觉空间定位跟踪（Inside-out 与 Outside-in 跟踪）；震动与触觉力反馈系统。
    \end{itemize}
    \item \textbf{游戏平台的社会、法律、伦理与无障碍考量}：
    \begin{itemize}
        \item \textbf{可用性工程与人机工效}：玩家交互直觉、认知负荷控制、游戏人机交互启发式评估方法。
        \item \textbf{无障碍支持与包容性设计}：色盲色弱显示增强滤镜、字幕与屏幕朗读辅助、全键位无障碍重映射适配、相关法规遵从性评估（如 CVAA 法案）。
        \item \textbf{绿色计算与软硬件可持续性}：低碳渲染模式设计、平台功耗限制、设备生命周期环保回收与反计划报废设计。
    \end{itemize}
    \item \textbf{实时物理仿真与现代图形渲染系统}：
    \begin{itemize}
        \item \textbf{现代 CPU 与 GPU 体系架构}：费林分类法（Flynn's Taxonomy，SIMD/SIMT 并行架构）；GPU 图形计算阵列、显存控制器、PCIe 显卡总线、显示扫描输出引擎。
        \item \textbf{渲染流水线与硬件光栅化架构}：移动端分块延迟渲染 (TBDR)与桌面端立即渲染模式 (IMR)的内存带宽开销权衡。
        \item \textbf{核心算法、空间数据结构与数学模型}：确定性游戏主循环；空间加速层次树划分（八叉树 Octree、包围盒层次 BVH、KD-Tree）；视锥体裁剪与细节层次（LOD）动态调度；快速碰撞检测（分离轴定理 SAT、GJK 算法）与刚体碰撞响应；NPC 行为树与分层有限状态机（HFSM）；程序化内容生成（PCG）。
        \item \textbf{多媒体虚拟世界资产计算与表达}：空间多声道三维音频、精灵图集、3D 网格模型与物理基础材质（PBR 材质双向反射分布函数 BRDF）、骨骼动画与逆运动学插值。
    \end{itemize}
    \item \textbf{游戏工程研发技术栈与主流工具链}：
    \begin{itemize}
        \item \textbf{游戏编程语言体系}：C++（高性能引擎核心）、C\#（Gameplay 逻辑）、Lua/Python（轻量热更新脚本）。
        \item \textbf{现代着色器语言与可视化着色器}：HLSL、GLSL、MSL 及基于节点的 ShaderGraph 可视化渲染材质系统。
        \item \textbf{底层现代图形 API}：DirectX 12、Vulkan、Metal、WebGL / WebGPU 底层显存显卡控制。
        \item \textbf{主流商业与开源游戏引擎}：Unreal Engine、Unity、Godot、CryEngine 等现代引擎架构；版本控制工具（处理海量二进制资产的 Git LFS、Perforce）；GPU 性能剖析分析器（RenderDoc、PIX、NVIDIA Nsight）。
        \item \textbf{关键开发工程方法}：快速原型验证、可执行构建流水线编译、真机帧率与显存性能调优、跨平台架构抽象移植、网络联机架构（状态同步、帧同步与预测回滚 Rollback 算法）。
    \end{itemize}
    \item \textbf{现代游戏设计工程体系与泛在应用}：
    \begin{itemize}
        \item \textbf{设计语言与模型}：MDA 理论框架（机制 Mechanics、动态 Dynamics、美感 Aesthetics）、心流体验模型。
        \item \textbf{以玩家为中心的设计思维}：双钻设计模型、迭代原型开发、玩家社区参与。
        \item \textbf{核心流派架构特征}：角色扮演（RPG）、第一人称射击（FPS）、即时战略（RTS）、多人在线战术竞技（MOBA）。
        \item \textbf{受众细分与玩家行为学}：巴特尔玩家分类模型（成就者、探索者、社交者、杀手）、游戏动机量表。
        \item \textbf{游戏技术的跨行业溢出与泛在应用}：严肃游戏、工业数字孪生、虚拟制片、全沉浸式模拟器与医学仿真培训。
    \end{itemize}
\end{itemize}

\textbf{【示范性学习成果与能力指标】}
\begin{itemize}
    \item 阐述现代 GPGPU 架构的核心特征，说明 SIMT 线程分支发散对游戏渲染效率的负面影响。
    \item 梳理图形光栅化渲染流水线从顶点处理、图元装配、光栅化、片元着色到模板/深度测试与输出合并的完整阶段。
    \item 运用现代着色器语言，按照光照反射模型要求手写实现一个自定义材质片元着色器。
    \item 利用现代图形 API 或工业级引擎，独立实现一个具备物理碰撞检测、空间四叉树分区与相机漫游的 3D 虚拟场景。
    \item 运用 GPU Profiler 性能分析工具，对一款游戏原型进行每帧 CPU 耗时、Draw Call 数量、显存带宽及过度绘制的专项诊断与优化。
    \item 针对跨平台游戏发布场景，评估一套能兼顾桌面 PC 与便携移动终端的高可移植性软件架构设计方案。
\end{itemize}

\subsection{交互式计算平台与大模型应用}
\textbf{【学时建议】}：2 领域核心学时（领域核心）。

本单元探讨支持人类进行探索性数据分析、实时可视化交互、创意编程，以及结合现代大语言模型（LLMs，如 OpenAI ChatGPT、Codex 与 GitHub Copilot）构建新型人机协同智能应用的前沿软件平台与范式。

\textbf{【核心知识点与专题大纲】}
\begin{itemize}
    \item \textbf{交互式数据探索分析计算平台}：
    \begin{itemize}
        \item 文档型可执行笔记本环境：Jupyter Notebook、Google Colab、Observable HQ、RStudio。
        \item 云端大数据交互查询平台：云 SQL 交互引擎（Google BigQuery、Snowflake）与分布式内存分析平台。
    \end{itemize}
    \item \textbf{交互式数据可视化生态与低时延渲染反馈}：
    \begin{itemize}
        \item 数据驱动的交互图表展示体系（D3.js、ECharts、Plotly）。
        \item 极低响应延迟的交互反馈闭环：硬件加速图形渲染工具、交互式参数调节画布。
    \end{itemize}
    \item \textbf{创意编程与生成式艺术交互框架}：
    \begin{itemize}
        \item 实时交互式媒体与艺术计算框架（Processing、p5.js、TouchDesigner、Max/MSP）。
        \item 现场即兴现场编码（Live Coding，如 Sonic Pi、TidalCycles 现场音乐生成）。
        \item 人机交互协同生成（人工智能/ML 交互式配对创作体验）。
    \end{itemize}
    \item \textbf{大语言模型（LLM）开发生态与协同应用}：
    \begin{itemize}
        \item 主流大模型编程助手与 API 集成：OpenAI ChatGPT/Codex API、GitHub Copilot、Anthropic Claude、本地开源推理平台（Ollama、vLLM）。
        \item \textbf{辅助数理与工科计算研究}：音频信号时域与傅里叶频域分析交互、数据统计建模与算法原型敏捷验证。
        \item \textbf{赋能人文社科与数字创意探索}：基于数据的交互式新闻叙事、文学脚本生成与音乐和声理论探索。
    \end{itemize}
\end{itemize}

\textbf{【示范性学习成果与能力指标】}
\begin{itemize}
    \item 在 Jupyter Notebook 或 Google Colab 平台中交互式加载并分析大规模复杂数据集，完成数据清洗、多维统计与深度交互图表绘制。
    \item 基于实时现场编程环境（如 Sonic Pi），编写一段能接收用户参数输入并实时调制音频合成器的音乐演奏程序。
    \item 利用主流大模型 API 构建一个具备特定上下文提示工程与函数调用能力的智能助手程序。
    \item 对比分析 AI 辅助平台自动生成的代码与人类专家编写的代码在时间复杂度、边缘边界条件防御以及长期可维护性上的本质差异。
    \item 阐明交互式探索计算与严谨工业级生产数据管道在开发范式、错误处理与吞吐量设计上的核心分歧。
\end{itemize}

\subsection{平台社会责任、伦理合规与专业素养}
\textbf{【学时建议】}：3 领域核心学时（领域核心）。

本单元立足专用平台开发的工程前沿，系统探讨移动端传感数据合规、物联网端侧监控、沉浸式元宇宙伦理、人机协作机器人生理心理安全以及界面多元化包容性（DEI）等重大社会与伦理议题。

\textbf{【核心知识点与专题大纲】}
\begin{itemize}
    \item \textbf{沉浸式扩展现实与增强技术的社会伦理冲击}：VR/AR 头戴显示器沉浸式体验对人类认知模式与现实脱节风险；元宇宙虚拟经济中的知识产权与未成年人心理防护。
    \item \textbf{具身机器人与自动化系统的人机伦理困境}：自动驾驶汽车事故责任归属与“电车难题”抉择；工业协作机器人的人身物理安全防护边界；服务型机器人拟人化带来的隐私边界模糊。
    \item \textbf{多元化、公平性与包容性（DEI）在界面设计中的工程落地}：适老化界面改造、消除对少数群体与残障群体的算法偏见；多语言全球化文化适配与色彩无歧视呈现。
    \item \textbf{端侧设备数据隐私法规与监控防范}：移动定位信息滥用规避、智能家居物联网麦克风与摄像头端侧隐私保护；GDPR 及相关个人信息保护法规在边缘计算设备上的合规实施。
\end{itemize}

\section{专业素养与思维倾向}

专用平台开发面向高度碎片化、快速迭代且异构的软硬件生态，要求计算机专业学生具备以下核心专业素养：
\begin{itemize}
    \item \textbf{自主学习与前沿技术追踪力}：新型硬件架构、专用操作系统与领域专用语言（DSL）层出不穷，工程师必须具备迅速阅读技术规范、脱离既有思维定势快速掌握全新平台开发工具链的元学习能力。
    \item \textbf{严苛约束下的工程创造力}：在电池功耗、内存空间、计算散热及通信带宽受到严厉约束的非常规平台上，能够创造性地设计高效算法与精简数据结构，达成极致性能优化。
    \item \textbf{异构环境敏捷适应性}：能够灵活游刃于完全不同的软硬件生态之间，面对移动端碎机型、嵌入式总线偶发噪声及云端动态漂移环境迅速调整工程实现方案。
    \item \textbf{严谨求实的软硬件联调素养 (Meticulousness in Debugging \& Testing)}：在面对跨时钟域总线、异步网络波动以及多线程并发竞态条件时，保持极端细致的逻辑严密性，善用硬件示波器、底层 Profiler 及系统日志深挖隐蔽 Bug。
\end{itemize}

\section{数学与跨学科基础要求}

专用平台软件开发横跨信号处理、机械运动控制、三维渲染管线与空间几何计算，对以下数学领域提出了明确要求：
\begin{itemize}
    \item \textbf{微积分}：用于机器人动力学中质点加速度与力矩的微分方程推导、物理引擎中位移积分计算以及连续物理信号采集滤波。
    \item \textbf{线性代数}：用于三维空间坐标系旋转平移（齐次坐标矩阵、四元数姿态插值变换）、机器人正逆运动学雅可比矩阵计算，以及现代 GPU 顶点几何渲染流水线变换。
    \item \textbf{概率论与数理统计}：用于物联网传感器噪声过滤（卡尔曼滤波算法）、弱网通信抖动与丢包率概率模型构建、大规模交互用户行为特征分布分析。
    \item \textbf{离散数学与结构}：用于自主导航算法中的图论寻路（$A^*$、Dijkstra）、机器人状态机与游戏行为树模型构建、事件驱动系统优先队列设计。
\end{itemize}

\section{交叉主题与跨学科共享概念}

专用平台开发是多门经典计算机学科的集大成者，与以下知识领域高度交叉：
\begin{itemize}
    \item \textbf{人工智能 (AI)}：端侧轻量化神经网络模型推理（TinyML）、大模型 API 赋能平台交互、机器人环境感知与自主决策规划。
    \item \textbf{计算机图形学 (GIT)}：三维网格几何表达、实时光栅化渲染、自定义片元/顶点着色器编写、虚拟现实沉浸式渲染。
    \item \textbf{人机交互 (HCI)}：触觉与多模态传感交互、以人为本的界面工效学设计、WCAG 无障碍可访问性开发标准。
    \item \textbf{系统建模与仿真}：刚体动力学仿真、高保真虚拟物理环境构建、数字孪生技术。
    \item \textbf{编程语言理论 (FPL)}：特定平台专用领域语言（DSL）、着色器着色语言机制、元编程与跨平台运行时的解释与 JIT 编译。
    \item \textbf{软件工程 (SE)}：模块解耦设计、持续集成与自动化发布流水线（CI/CD）、跨平台代码复用度管理、微服务治理。
    \item \textbf{需求工程}：深入挖掘垂直领域用户特殊场景痛点，在设备物理制约与业务功能体验之间达成最佳工程平衡。
\end{itemize}

\section{课程体系规划与开课建议}

各高校可依据自身师资力量与专业特色方向，将本领域内容规划组合为以下典型专业模块课程：
\begin{enumerate}
    \item \textbf{移动与现代 Web 应用全栈开发 (Mobile and Modern Web App Development, 48 学时)}：
    全面覆盖 Web 前端组件化与云端微服务基础（24h）、移动原生与跨平台架构实战（18h）、数据安全通信与端侧轻量数据库（6h）。侧重从零搭建一套具备云端弹性后端支撑的高性能多端移动商业应用。
    \item \textbf{机器人软件工程与 ROS 系统实战 (Robotics Software Engineering and ROS, 45 学时)}：
    深入讲解 Linux 实时环境与 C++ 开发（9h）、ROS2 节点/话题/服务架构与 Gazebo 仿真（15h）、机器人运动学与 SLAM 自主导航建图（15h）、实体差速底盘或机械臂硬件抓取实验（6h）。
    \item \textbf{嵌入式系统与物联网全栈开发 (Embedded Systems and IoT Development, 45 学时)}：
    涵盖微控制器体系与 C 语言裸机外设编程（12h）、FreeRTOS 任务调度与多任务并发同步（15h）、低功耗无线物联网协议接入（BLE/Wi-Fi/MQTT, 12h）、硬件在环故障诊断实训（6h）。
    \item \textbf{游戏引擎架构与实时交互图形系统 (Game Engine Architecture and Interactive Media, 48$\sim$60 学时)}：
    涵盖现代 GPU 架构与实时着色器编程（18h）、游戏主循环与空间碰撞加速结构（15h）、现代商业引擎核心架构与网络同步对战机制（15h）。
    \item \textbf{交互式数据计算与生成式 AI 编程 (Interactive Data Computing and LLM Programming, 32 学时)}：
    涵盖交互式 Notebook 数据探索与流式可视化（12h）、大模型 API 提示词工程与函数调用架构（12h）、基于 LLM 的智能助手与创意代码生成实战（8h）。
\end{enumerate}

\section{真实能力任务规范}

\begin{taskbox}[任务 1：原生开发与跨平台技术选型决策评估]
\textbf{【任务情境】}：某初创科技企业拟推出一款融合高帧率动画、复杂底层传感器调用及多端快速部署特征的数字化健康服务产品，技术团队在“全原生开发（iOS Swift + Android Kotlin）”与“跨平台框架”之间存在选型分歧。\\
\textbf{【能力陈述】}：综合技术架构与产品设计维度，权衡系统运行性能、包体积、原生硬件接口访问能力、发布周期及人力维护成本，撰写客观详实的架构技术选型论证报告。\\
\textbf{【能力领域】}：应用开发 \quad \textbf{【能力单元】}：系统评估\\
\textbf{【关联知识单元】}：软件工程/工程工具与开发环境、专用平台/专用平台共性特征、专用平台/移动平台\\
\textbf{【达标等级】}：\textbf{解释与评估} \quad \textbf{【学时层级】}：计算机核心
\end{taskbox}

\begin{taskbox}[任务 2：高一致性多端自适应移动应用设计与实现]
\textbf{【任务情境】}：企业核心业务移动客户端需要在各厂商碎片化 Android 设备、多样化分辨率折叠屏以及不同版本 iOS 系统上保持高度一致的用户体验与流畅动画。\\
\textbf{【能力陈述】}：熟练运用现代移动开发语言与 UI 响应式布局体系，优化应用触控响应与内存占用，通过真机云测与自动化 UI 测试确保在跨机型、异构操作系统上稳定达标 60 FPS 刷新率。\\
\textbf{【能力领域】}：软件与应用 \quad \textbf{【能力单元】}：设计、开发与测试\\
\textbf{【关联知识单元】}：人机交互/理解用户、编程语言/面向对象编程、软件开发基础/开发方法与实践、软件工程/软件构造、软件工程/测试与质量保证、专用平台/移动平台\\
\textbf{【达标等级】}：\textbf{开发} \quad \textbf{【学时层级】}：计算机核心
\end{taskbox}

\begin{taskbox}[任务 3：基于引擎将策划蓝图原型重构为高性能跨平台核心系统]
\textbf{【任务情境】}：游戏策划与交互设计师在商业游戏引擎中利用可视化脚本（如 Unreal Blueprints）快速搭建了玩法原型，但在高密度复杂场景下出现严重的 CPU 主线程掉帧，且需移植发布至 PC、游戏主机与移动终端。\\
\textbf{【能力陈述】}：熟练剖析游戏主循环与底层渲染瓶颈，将可视化原型逻辑重构为基于 C++ 的高性能模块，运用空间数据结构与多线程机制优化碰撞计算，确保核心玩法的跨平台无缝兼容与高流畅度。\\
\textbf{【能力领域】}：软件与系统应用 \quad \textbf{【能力单元】}：系统研发、性能评测与调优\\
\textbf{【关联知识单元】}：软件工程/软件架构与设计、软件工程/软件构造、算法/算法策略与范式、系统基础/系统性能、专用平台/专用平台共性特征、专用平台/游戏平台\\
\textbf{【达标等级】}：\textbf{开发} \quad \textbf{【学时层级】}：领域核心
\end{taskbox}

\begin{taskbox}[任务 4：沉浸式头显（HMD）与混合现实新型外设交互可用性与无障碍评估]
\textbf{【任务情境】}：某团队研发了一款配合空间计算头戴式显示器（HMD）使用的手部骨骼力反馈外设，拟面向公共教育与游戏互动场景进行广泛推广。\\
\textbf{【能力陈述】}：制定严谨的用户人机工效测试方案，针对不同手型、行动障碍以及弱视用户群体开展可用性与无障碍合规性评估，针对前庭动晕症（晕动症）与力反馈延迟提出工程改进建议。\\
\textbf{【能力领域】}：系统与人机体验 \quad \textbf{【能力单元】}：需求分析、用户测试与社会适应性\\
\textbf{【关联知识单元】}：人机交互/理解用户、人机交互/无障碍与包容性设计、图形交互/沉浸式技术、图形交互/物理计算、社会伦理/多样性、平等与包容 (DEI)、专用平台/游戏平台\\
\textbf{【达标等级】}：\textbf{应用与评估} \quad \textbf{【学时层级】}：领域核心
\end{taskbox}

\begin{taskbox}[任务 5：基于演进商业需求构建高安全与高弹性 Web 业务系统]
\textbf{【任务情境】}：某企业数字化交易平台随着业务扩张面临跨国数据合规审计与日益频繁的 Web 漏洞攻击风险，亟需重构底层业务架构。\\
\textbf{【能力陈述】}：熟练运用全栈现代编程语言与云服务组件，实施防范 XSS、CSRF 与 SQL 注入的纵深安全防御，构建多可用区弹性微服务底座，确保系统在流量洪峰下兼顾安全与高可用。\\
\textbf{【能力领域】}：应用工程 \quad \textbf{【能力单元】}：系统设计与安全评估\\
\textbf{【关联知识单元】}：体系结构/性能与能效、网络通信/网络空间安全、操作系统/安全防护与保护机制、系统基础/系统安全、软件工程/软件架构与设计、专用平台/专用平台共性特征、专用平台/Web 平台、社会伦理/隐私保护\\
\textbf{【达标等级】}：\textbf{开发} \quad \textbf{【学时层级】}：计算机核心
\end{taskbox}

\begin{taskbox}[任务 6：面向移动计算生态设计与交付高性能移动业务 API]
\textbf{【任务情境】}：公司需要重构移动应用数据层交互通道，以往臃肿的桌面端通用 REST 接口在移动弱网环境下导致严重的客户端请求堆叠与电池电量过耗。\\
\textbf{【能力陈述】}：遵循移动平台最佳工程规范，设计轻量级、低冗余、支持批量聚合与断点续传的移动专属 API 接口规范，编写高可靠移动端数据网络层 SDK 并实现完善的本地离线缓存支持。\\
\textbf{【能力领域】}：应用软件工程 \quad \textbf{【能力单元】}：接口研发与性能评测\\
\textbf{【关联知识单元】}：编程语言/硬件接口、操作系统/文件系统 API 与底层实现、软件工程/软件架构与设计、软件工程/工程工具与开发环境、专用平台/专用平台共性特征、专用平台/移动平台\\
\textbf{【达标等级】}：\textbf{开发} \quad \textbf{【学时层级】}：计算机核心
\end{taskbox}

\begin{taskbox}[任务 7：自主智能机器人软硬件交互系统端到端设计与实现]
\textbf{【任务情境】}：针对自动化工厂物流中周转物料的自动寻引与精确装卸需求，团队需要开发一款具备自主避障定位与末端抓取的机器人系统。\\
\textbf{【能力陈述】}：准确梳理工业场景核心任务需求，设计基于 ROS2 架构的软件节点通信拓扑，打通激光雷达建图、底盘逆运动学解算与末端机械臂动作控制，实现机器人全自动化运行。\\
\textbf{【能力领域】}：系统应用工程 \quad \textbf{【能力单元】}：系统设计、集成与评估\\
\textbf{【关联知识单元】}：算法/算法策略与范式、软件工程/工程工具与开发环境、软件工程/软件架构与设计、专用平台/专用平台共性特征、专用平台/机器人平台\\
\textbf{【达标等级】}：\textbf{解释与应用} \quad \textbf{【学时层级】}：领域核心
\end{taskbox}

\begin{taskbox}[任务 8：面向严苛硬件指标的高可靠嵌入式系统架构设计与协议研发]
\textbf{【任务情境】}：医疗设备研发团队需开发一款便携式多参数生理监护仪，要求系统在毫瓦级低功耗下持续采样心电信号，且软硬件系统绝不允许发生阻塞死锁。\\
\textbf{【能力陈述】}：依据硬件芯片物理规格书，在轻量实时操作系统（RTOS）上设计多任务确定性调度体系，实现 SPI 传感器总线高频通信与看门狗自动故障恢复机制，编写完整的板级支持包（BSP）。\\
\textbf{【能力领域】}：系统应用工程 \quad \textbf{【能力单元】}：系统架构研发、软硬件部署与集成\\
\textbf{【关联知识单元】}：编程语言/硬件接口、软件工程/软件架构与设计、软件工程/工程工具与开发环境、系统基础/系统性能、专用平台/专用平台共性特征、专用平台/嵌入式平台\\
\textbf{【达标等级】}：\textbf{开发} \quad \textbf{【学时层级】}：领域核心
\end{taskbox}

\begin{taskbox}[任务 9：大规模分布式 Web 站点与线上资产持续运维与性能优化]
\textbf{【任务情境】}：某跨国电商企业在“大促”活动前夕，线上多套 Web 站点与核心交易网关面临瞬时百万级并发写入，底层数据库与反向代理服务器偶发超时告警。\\
\textbf{【能力陈述】}：深入分析系统网络吞吐、单跳链路与 NoSQL 分布式存储集群瓶颈，部署 CDN 静态资产边缘缓存、微前端动态路由分发与读写分离集群，建立自动化监控与灾备熔断机制。\\
\textbf{【能力领域】}：系统运维与工程应用 \quad \textbf{【能力单元】}：系统集成、监控与性能调优\\
\textbf{【关联知识单元】}：数据管理/NoSQL 系统与非关系型存储、软件工程/工程工具与开发环境、软件工程/软件架构与设计、网络通信/单跳通信与局域网、操作系统/文件系统 API 与底层实现、专用平台/专用平台共性特征、专用平台/Web 平台\\
\textbf{【达标等级】}：\textbf{应用与调优} \quad \textbf{【学时层级】}：计算机核心
\end{taskbox}

\begin{taskbox}[任务 10：跨职能团队协同交付企业级高可用全栈平台解决方案]
\textbf{【任务情境】}：在现代数字化敏捷研发流程中，涉及前端工程师、后端微服务架构师、移动客户端工程师、UI/UX 设计师以及产品运营经理等多角色的紧密协作。\\
\textbf{【能力陈述】}：展现良好的团队协作与技术沟通能力，综合梳理前后端交互协议、数据实体约束与客户端交互边界，制定标准化 API 契约与持续集成（CI）流水线，高质量按期交付全栈产品。\\
\textbf{【能力领域】}：工程管理与软件系统 \quad \textbf{【能力单元】}：团队协作、部署与交付\\
\textbf{【关联知识单元】}：软件工程/项目管理与度量、软件工程/软件架构与设计、软件工程/团队协作、专用平台/专用平台共性特征、专用平台/Web 平台、专用平台/移动平台\\
\textbf{【达标等级】}：\textbf{解释与协作} \quad \textbf{【学时层级】}：领域通用
\end{taskbox}

\section{指导委员会与主要贡献者}

\begin{itemize}
    \item \textbf{工作组主席}：
    \begin{itemize}
        \item Christian Servin (El Paso Community College, 美国)
    \end{itemize}
    \item \textbf{委员会核心委员}：
    \begin{itemize}
        \item Sherif G. Aly (The American University in Cairo, 埃及)
        \item Yoonsik Cheon (The University of Texas at El Paso, 美国)
        \item Eric Eaton (University of Pennsylvania, 美国)
        \item Claudia L. Guevara (Jochen Schweizer mydays Holding GmbH, 德国)
        \item Larry Heimann (Carnegie Mellon University, 美国)
        \item Amruth N. Kumar (Ramapo College of New Jersey, 美国)
        \item R. Tyler Pirtle (Google, 美国)
        \item Michael Scott (Falmouth University, 英国)
    \end{itemize}
    \item \textbf{主要贡献学者与行业专家}：
    \begin{itemize}
        \item Orlando Gordillo (NASA, 美国)
        \item Sean R. Piotrowski (Rider University, 美国)
        \item Mark O'Neil (Blackboard Inc., 美国)
        \item John DiGennaro (Qwickly, 美国)
        \item Rory K. Summerley (Falmouth University, 英国)
    \end{itemize}
\end{itemize}

\begin{strip}
\section{附录：核心主题与技能要求等级对照表}

\begingroup
\zihao{-5}\renewcommand{\arraystretch}{0.90}\setlength{\tabcolsep}{3pt}
\noindent\begin{tabularx}{\textwidth}{p{0.17\textwidth} Y c c c}
\toprule
\textbf{知识单元} & \textbf{核心知识点与具体主题} & \textbf{认知等级} & \textbf{分类} & \textbf{学时} \\
\midrule
\multirow{4}{*}{\makecell[l]{专用平台/跨平台\\共性技术}} 
& 专用平台生态全景（Web、移动端、机器人、嵌入式、游戏平台） & 解释 & 计算机核心 & 0.5 \\
& 专用平台 API 编程契约 vs. 传统桌面通用软件构建 & 应用 & 计算机核心 & 0.5 \\
& 跨平台开发语言图谱（Kotlin、Swift、C\#、C++、TypeScript 等） & 应用 & 计算机核心 & 0.5 \\
& 严苛资源受限（算力、内存、通信、功耗、可用性）下的工程构建 & 应用 & 计算机核心 & 0.5 \\
\midrule
\multirow{4}{*}{\makecell[l]{专用平台/Web 平台\\开发体系}} 
& 现代 Web 编程语言标准与前端组件化元框架 & 应用 & 计算机核心 & 1 \\
& SaaS 多租户架构与 W3C/WCAG 无障碍合规性规范 & 解释 & 计算机核心 & 0.5 \\
& Web 纵深安全防御机制（XSS、CSRF、CSP 与跨域通信） & 评估 & 计算机核心 & 0.5 \\
& 云原生弹性伸缩、无服务器与关系/NoSQL 存储选型 & 应用 & 计算机核心 & 1 \\
\midrule
\multirow{4}{*}{\makecell[l]{专用平台/移动计算\\平台开发}} 
& 移动端原生开发与跨平台开发框架架构选型对比 & 评估 & 计算机核心 & 0.5 \\
& 移动平台硬性软硬件约束（触控设计、功耗与性能权衡、沙箱安全） & 解释 & 计算机核心 & 0.5 \\
& 移动端特有使能技术（LBS 位置感知、多传感器融合、AR 空间计算） & 应用 & 计算机核心 & 1 \\
& 现代异步非阻塞编程范式（协程、async/await）与弱网容错测试 & 应用 & 计算机核心 & 1 \\
\midrule
\multirow{3}{*}{\makecell[l]{专用平台/机器人\\开发平台}} 
& 机器人平台类型、板载嵌入式计算、传感器与伺服驱动执行器 & 解释 & 领域核心 & 1 \\
& 机器人操作系统 ROS/ROS2 体系架构与软硬件虚实仿真互联 & 应用 & 领域核心 & 1 \\
& 正逆运动学、刚体动力学基础、SLAM 室内建图与自主路径规划 & 应用 & 领域核心 & 1 \\
\midrule
\multirow{3}{*}{\makecell[l]{专用平台/嵌入式与\\物联网平台}} 
& 嵌入式微控制器特性、硬实时/软实时时延确定性与功能安全分析 & 解释 & 领域核心 & 1 \\
& 物理感知外设硬件接口驱动（GPIO、ADC/DAC、SPI、I2C、CAN） & 应用 & 领域核心 & 1 \\
& 实时操作系统 RTOS 任务调度、低功耗休眠管理与硬件在环调试 & 应用 & 领域核心 & 1 \\
\midrule
\multirow{4}{*}{\makecell[l]{专用平台/游戏引擎\\互动平台}} 
& 游戏硬件演进、显示与外设跟踪接口、无障碍与绿色计算规范 & 解释 & 领域核心 & 1 \\
& GPU 现代并行架构、TBDR 与立即渲染流水线对比、着色器编程 & 应用 & 领域核心 & 1 \\
& 实时仿真算法（游戏主循环、八叉树/BVH 空间划分、刚体碰撞检测） & 应用 & 领域核心 & 1 \\
& 商业游戏引擎架构、跨平台移植调优与多端网络同步预测回滚 & 应用 & 领域核心 & 1 \\
\midrule
\multirow{2}{*}{\makecell[l]{专用平台/交互计算\\与大模型平台}} 
& 交互式数据分析探索平台与低时延动态可视化 & 应用 & 领域核心 & 1 \\
& 大语言模型（LLMs）开发生态、API 集成与人机协同编程实践 & 应用 & 领域核心 & 1 \\
\midrule
\makecell[l]{专用平台/社会伦理\\与专业素养} 
& 沉浸式技术现实冲击、机器人伦理安全、适老化 DEI 与端侧隐私合规 & 评估 & 领域核心 & 3 \\
\bottomrule
\end{tabularx}
\endgroup
\end{strip}
